% !TeX root = main.tex
% Build with latexmk so BibTeX runs before the final LaTeX passes.
\documentclass[10pt,conference]{IEEEtran}

\usepackage{amsmath}
\usepackage{booktabs}
\usepackage{cite}
\usepackage[hidelinks]{hyperref}
\usepackage{microtype}

\title{Fairness in Adaptive Quantum-Network Routing:\\
Service-Equity Risks and Mitigations for Bandit-Based Qubit Allocation}

\author{\IEEEauthorblockN{Piter Z. Garcia Bautista}
\IEEEauthorblockA{DSCI 602: Applied Data Science II\\
Capstone: Fairness-Aware Adaptive Routing for Quantum Networks\\
Chosen Focus Area: Fairness}}

\begin{document}
\maketitle

\begin{abstract}
This analysis examines fairness risks in an adaptive quantum-network controller that uses bandit policies to select entanglement routes and allocate limited qubit-generation attempts. Aggregate throughput can conceal systematic service degradation for flows with higher demand, weaker context, or less favorable topology. The appropriate objective is therefore not classification parity alone, but demand-normalized service equity: comparable access to useful quantum service after accounting for legitimate demand and physical constraints. The analysis identifies pipeline and feedback-loop harms, specifies measurable service-equity criteria, and proposes constrained mediation, stratified stress testing, and auditable decision traces. These mitigations reduce but cannot eliminate the tension among fairness, fidelity, latency, robustness, and total network utility.
\end{abstract}

\section{Context and Fairness Objective}

The capstone extends an existing threat-aware quantum-routing framework with fairness-aware mediation. At each decision round, a controller observes a possibly incomplete network state, selects a route, and allocates scarce qubit or entanglement-generation capacity. The system influences whether competing requests receive successful end-to-end entanglement, how long they wait, how many retries they incur, and whether repeated allocation decisions starve a flow. The directly affected stakeholders are users and applications submitting concurrent requests; network operators are accountable for how the controller translates declared demand and service requirements into access to quantum resources.

The relevant groups are operational service classes or flow cohorts defined before evaluation, such as latency-sensitive versus delay-tolerant traffic, high- versus low-demand requests, or flows exposed to high- versus low-quality context. These groups must not be invented after observing results. If future deployments associate service classes with people or institutions, protected-attribute analysis would also be required, but this paper does not assume demographic data that the current simulator does not contain.

Classical group criteria provide useful warnings but are not direct substitutes for network-allocation fairness. Demographic parity can be inappropriate when legitimate demands differ, while equalized-odds and calibration criteria concern prediction errors and can conflict when base rates differ~\cite{hardt2016equality,chouldechova2017fair}. The primary criterion here is therefore \emph{demand-normalized service equity}, supported by outcome gaps and a starvation constraint. For group $g$ over horizon $T$, define normalized access
\begin{equation}
 A_g(T)=\frac{\sum_{t=1}^{T}Q_{g,t}}{\sum_{t=1}^{T}D_{g,t}},
 \qquad
 \Delta_A(T)=\max_{g,h}|A_g(T)-A_h(T)|,
\end{equation}
where $Q_{g,t}$ is allocated quantum capacity and $D_{g,t}$ is declared demand. The audit also reports entanglement-success and latency gaps, $\Delta_S=\max_{g,h}|P(\mathrm{success}\mid g)-P(\mathrm{success}\mid h)|$ and $\Delta_L=\max_{g,h}|\bar L_g-\bar L_h|$, plus Jain's index over group-normalized access~\cite{jain1984quantitative}. This combination measures allocation, experienced service, and distribution rather than declaring one scalar to be ``fairness.'' Quantum-routing research already shows that finite capacity must be scheduled across simultaneous requests and that progressive filling can approach max--min fairness~\cite{li2021effective}; fair-sharing work in distributed quantum computing further establishes service differentiation as a quantum-resource problem~\cite{cicconetti2023service}.

\section{Risk Analysis}

Fairness failures can enter at several points. Suresh and Guttag distinguish historical, representation, measurement, aggregation, learning, evaluation, and deployment harms across the machine-learning life cycle~\cite{suresh2021framework}; broader surveys similarly warn that bias is produced by interacting data, model, and institutional choices rather than by data alone~\cite{mehrabi2021survey}. In this capstone, \emph{measurement bias} arises when link quality, fidelity, queue state, or demand is measured less accurately for one cohort. \emph{representation bias} arises when stress tests omit difficult topologies, service classes, or context-quality regimes. \emph{aggregation bias} arises when a single policy assumes that flows with different latency or fidelity requirements should share one reward model. \emph{evaluation bias} occurs when aggregate reward or throughput masks subgroup starvation and tail latency.

Sequential learning adds a feedback risk. A bandit observes rewards only for selected actions; underexplored routes or cohorts generate less evidence, which can lower confidence and reduce future selection. Joseph et al. show that fairness constraints change the regret cost of learning under uncertainty~\cite{joseph2016fairness}. Chen et al. demonstrate that contextual allocation can impose minimum service rates and that context matters when users perform differently across tasks~\cite{chen2020fair}. For quantum routing, stale or systematically weaker context can therefore become an allocation disadvantage even when the controller never uses a protected attribute. Threats and nonstationarity can intensify the problem: a policy may preserve acceptable global efficiency while adaptive disruptions concentrate failures, retries, or delay on a minority service class.

The plausible harms are concrete. A latency-sensitive class may repeatedly receive routes with longer queues; high-demand flows may receive equal absolute allocations but lower demand-normalized access; poorly observed flows may be treated as low-value; and an efficiency-maximizing allocator may repeatedly favor short or reliable routes until harder requests starve. These are allocative harms because access to a scarce network capability is withheld or degraded. They also create reliability and accountability harms when operators cannot reconstruct why a request received inferior service.

\section{Mitigations and Validation}

First, the data and trace contract should record, for every decision, the service-group definition, declared demand, context-quality indicators, candidate routes, policy scores, selected route, allocated qubits, fidelity estimate, success, latency, retry count, and starvation age. Group definitions, exclusions, and physical infeasibility conditions should be versioned with each run. This makes disparity calculations reproducible and prevents post-hoc redefinition.

Second, fairness should enter the decision loop through a bounded mediator rather than replace the base routing policy. A candidate score is
\begin{equation}
 S_t^{\mathrm{fair}}(a)=U_t(a)-\lambda[\widehat{\Delta}_t(a)-\epsilon]_+-\gamma C_t(a),
\end{equation}
where $U_t(a)$ is the underlying utility score, $\widehat{\Delta}_t(a)$ estimates the action's disparity effect, $\epsilon$ is an allowed gap, and $C_t(a)$ penalizes missing, stale, or low-confidence context. A minimum-service or maximum-starvation constraint should remain active even when the soft penalty is insufficient. This design follows fair contextual allocation work while keeping the fairness intervention separable and auditable~\cite{chen2020fair}.

Third, validation must compare an unmediated baseline ($\lambda=0$) against multiple $\lambda$, $\gamma$, and $\epsilon$ settings across matched policies, allocators, capacities, topologies, threat regimes, and random seeds. Reports should include aggregate utility, regret, success, and fidelity alongside $\Delta_A$, $\Delta_S$, $\Delta_L$, Jain's index, worst-group outcomes, and starvation frequency. Confidence intervals and repeated trials are necessary because small cohorts can produce unstable gaps. A mitigation passes only if it improves prespecified equity measures without violating minimum fidelity, safety, or service constraints; results should expose the full utility--equity frontier rather than select a favorable operating point after inspection.

Finally, operators should use monitoring thresholds, periodic subgroup audits, drift tests, and an override or rollback path. When a group lacks enough observations for a stable estimate, the system should report uncertainty and apply a conservative service floor rather than infer that no disparity exists. The deployment documentation should identify who reviews alerts and who can change fairness weights.

\section{Limitations and Broader Impacts}

These mitigations cannot prove that a service taxonomy is socially legitimate, and equalizing measured outcomes may not correct structural differences in topology or access. Minimum allocations can reduce total throughput, increase latency for other flows, or spend qubits on routes that cannot meet fidelity requirements. Conversely, unconstrained efficiency can hide persistent exclusion. The acceptable trade-off is therefore a governance decision informed by measured frontiers, not a purely technical optimum. The present analysis is prospective: it proposes risks, metrics, and tests but does not claim that live mediation or DSCI602 fairness experiments are complete. Simulation evidence will also remain limited by modeled workloads and threats. A responsible deployment would require stakeholder review, hardware-grounded validation, transparent service-class rules, and continuing audits as network demand and topology change.

\clearpage
\bibliographystyle{IEEEtran}
\bibliography{fairness_references}

\end{document}
